
\documentclass[11pt]{article}

\usepackage[margin=1in]{geometry}
\usepackage[T1]{fontenc}
\usepackage[utf8]{inputenc}
\usepackage{lmodern}
\usepackage{microtype}
\usepackage{amsmath,amssymb,amsthm,mathtools}
\usepackage{booktabs,longtable,array}
\usepackage{enumitem}
\usepackage{hyperref}
\usepackage{xcolor}
\usepackage{url}

\hypersetup{
  colorlinks=true,
  linkcolor=blue!50!black,
  citecolor=blue!50!black,
  urlcolor=blue!50!black
}

\newtheorem*{theoremstar}{Theorem}
\newtheorem*{propositionstar}{Proposition}
\newtheorem*{corollarystar}{Corollary}
\newtheorem*{lemmastar}{Lemma}

\newcommand{\Omegaelem}{\Omega_{\mathrm{elem}}}
\newcommand{\Omegaprimeelem}{\Omega'_{\mathrm{elem}}}
\newcommand{\Omegafull}{\Omega_{\mathrm{full}}}
\newcommand{\CrossRules}{\operatorname{CrossRules}_{\sigma}}
\newcommand{\Assoc}{\operatorname{Assoc}}
\newcommand{\row}{\operatorname{row}}
\newcommand{\col}{\operatorname{col}}
\newcommand{\Fix}{\operatorname{Fix}}
\newcommand{\Aut}{\operatorname{Aut}}
\newcommand{\End}{\operatorname{End}}
\newcommand{\Abs}{\operatorname{Abs}}

\title{$\sigma$-Equivariant Elementary Magmas on $AG(2,3)$\\
\large Classification, associativity spectrum, and standard elementary selection in a six-rule family}
\author{
Burundai Taryi\\
Independent researcher\\
\texttt{burundai.taryi@gmail.com}
}
\date{April 21, 2026}

\begin{document}
\maketitle

\begin{abstract}

We study a finite elementary family of \(\sigma\)-equivariant row-graded magmas on
\(M=S\times S\), where \(|S|=3\). The complete space of \(\sigma\)-equivariant
cross-row rules has cardinality \(3^{18}\). The present paper restricts to six
distinguished elementary cross-rules \(g_1,\ldots,g_6\) and all \(27\)
\(\sigma\)-equivariant diagonal maps. The resulting elementary space
\[
\Omega_{\mathrm{elem}}=\{g_1,\ldots,g_6\}\times\Delta_\sigma
\]
has \(162\) objects and decomposes into \(90\) orbits under diagonal relabelling
by \(S_3\).

For this elementary family we derive an explicit master formula for the number
of associative triples in terms of the elementary cross-rule and diagonal
fixed-point data. Consequently, the associativity spectrum in
\(\Omega_{\mathrm{elem}}\) consists of exactly \(28\) values in \([162,567]\), all
divisible by \(3\). Among the six standard elementary magmas, two
independent criteria---minimum conditional column-access information and minimum
associativity---select the PAB magma. A restricted sequential selection theorem
further shows that PAB is uniquely selected inside
\[
\Omega'_{\mathrm{elem}}=\mathcal F_\sigma\times\{g_1,\ldots,g_6\}\times\Delta_\sigma
\]
by successively minimizing access information, diagonal entropy, and
associativity on the remaining candidate set. Finally, the PAB operation induces an
absorption graph on the six off-diagonal states generated by two transitions
\(T_\alpha,T_\beta\), naturally realizing the \(A_2\) root hexagon.

The full \(3^{18}\)-cross-rule universe, full isomorphism rigidity beyond the
diagonal \(S_3\)-action, operator-level path-memory invariants, and broader
Lie-theoretic interpretations remain outside the scope of this paper.

\end{abstract}


\section*{Scope and notation}
\addcontentsline{toc}{section}{Scope and notation}

All classification, orbit-count, master-formula, spectrum, and selection
statements below are scoped to one of the following spaces:
\[
\Omega_{\mathrm{elem}}=\{g_1,\ldots,g_6\}\times\Delta_\sigma,
\qquad |\Omega_{\mathrm{elem}}|=162,
\]
\[
\Omega'_{\mathrm{elem}}=\mathcal F_\sigma\times\{g_1,\ldots,g_6\}\times\Delta_\sigma,
\qquad |\Omega'_{\mathrm{elem}}|=1458,
\]
and the full cross-rule universe
\[
\Omega_{\mathrm{full}}=\operatorname{CrossRules}_\sigma\times\Delta_\sigma,
\qquad |\operatorname{CrossRules}_\sigma|=3^{18}.
\]
This paper classifies and analyses \(\Omega_{\mathrm{elem}}\), not
\(\Omega_{\mathrm{full}}\). The six-rule list is not exhaustive: it is a restricted
family of distinguished elementary formulas inside the full set of
\(3^{18}\) equivariant cross-row rules.



\section{Introduction}

\subsection{Object of study}

Let
\[
S=\{0,1,2\},\qquad M=S\times S.
\]
Elements of \(M\) are written \((r,c)\), with row \(r\) and column \(c\). The cyclic
shift \(\sigma:i\mapsto i+1\pmod 3\) acts diagonally on \(M\):
\[
\sigma(r,c)=(r+1,c+1).
\]

We study row-graded binary operations on \(M\), meaning
\[
(x\cdot y)_{\mathrm{row}}=x_{\mathrm{row}}.
\]
Inside each row we use the three-point Steiner rule, while self-products are
controlled by a \(\sigma\)-equivariant diagonal map. Cross-row multiplication is
the main classification variable.

A general \(\sigma\)-equivariant cross-row rule has \(3^{18}\) possibilities. We
isolate six elementary closed-form rules \(g_1,\ldots,g_6\). These six rules are
small enough to admit a complete associativity formula, spectrum computation,
and selection theory, while still containing the PAB magma. The general
background for finite non-associative systems, designs, finite geometries, and
root systems is standard; see \cite{bruck1958,denes1974,beth1999,colbourn2006,dembowski1968}.

\subsection{Main results}

\begin{theoremstar}[A --- PAB structure]
The PAB magma has exactly three idempotents,
automorphism group \(S_3\), endomorphism monoid
\(\{\widehat\rho:\rho\in S_3\}\cup\{\varepsilon_0,\varepsilon_1,\varepsilon_2\}\), and associativity count
\(219/729\). Its non-associativity is row-confined.

\end{theoremstar}
\begin{theoremstar}[B --- Elementary classification]
The full set of \(\sigma\)-equivariant
cross-row rules has cardinality \(3^{18}\). Restricting to six distinguished
elementary cross-rules and \(27\) diagonal maps gives
\(|\Omega_{\mathrm{elem}}|=162\). Under diagonal relabelling by \(S_3\), this space
has \(90\) orbits.

\end{theoremstar}
\begin{theoremstar}[C --- Associativity formula]
For every magma in
\(\Omega_{\mathrm{elem}}\), associativity is determined by the elementary
cross-rule and diagonal fixed-point data. The formula is verified over all
\(162\) magmas and all \(729\) triples.

\end{theoremstar}
\begin{theoremstar}[D --- Elementary spectrum]
The associativity spectrum of
\(\Omega_{\mathrm{elem}}\) consists of exactly \(28\) values in \([162,567]\), all
divisible by \(3\). The extrema \(162\) and \(567\) are attained.

\end{theoremstar}
\begin{theoremstar}[E --- Standard elementary selection]
Among the six standard
elementary magmas, \(g_1\) is the unique rule that simultaneously minimises
conditional column-access information and associativity. The selected magma is
PAB.

\end{theoremstar}
\begin{theoremstar}[F --- Restricted sequential selection]
Inside \(\Omega'_{\mathrm{elem}}\),
successive minimisation of \(H_{\mathrm{access}}\), then \(H_{\mathrm{diag}}\), and then
associativity on the remaining candidate set uniquely selects PAB.

\end{theoremstar}
\begin{theoremstar}[G --- Absorption graph]
For \(x\in M^\times\), the elements absorbed by
\(x\) are exactly \(T_\alpha(x)\) and \(T_\beta(x)\).

\end{theoremstar}
\subsection{Structure of the paper}

Section 2 defines PAB and proves its basic structure. Section 3 constructs the
elementary classification space and the diagonal \(S_3\)-orbit count. Section 4
derives the row-pattern associativity formula. Section 5 records the elementary
associativity spectrum. Section 6 proves the information-theoretic selection
results. Section 7 proves the absorption graph theorem. Appendix A gives the computational
verification artifacts supporting the finite claims.



\section{The PAB magma}

\subsection{Setup}

For each \(r\in S\), write
\[
F_r=\{(r,0),(r,1),(r,2)\}
\]
for the row fiber. For distinct \(i,j\in S\), write \(ij\) for the unique third
element of \(S\setminus\{i,j\}\).

\subsection{Definition}

The PAB magma is the algebra \((M,\cdot)\) defined as follows. Let
\(x=(r_1,c_1)\), \(y=(r_2,c_2)\).

\paragraph{AX0 --- row grading.} 
\[
(x\cdot y)_{\mathrm{row}}=r_1.
\]

\paragraph{AX1 --- cross-row rule.} If \(r_1\ne r_2\), then
\[
x\cdot y=(r_1,r_2).
\]

\paragraph{AX2 --- Steiner fiber rule.} If \(r_1=r_2\) and \(c_1\ne c_2\), then
\[
x\cdot y=(r_1,c_1c_2).
\]

\paragraph{AX3 --- diagonal collapse.} If \(r_1=r_2\) and \(c_1=c_2=c\), then
\[
(r_1,c)\cdot(r_1,c)=(r_1,r_1).
\]

\subsection{Idempotents and square fibers}

The idempotents are exactly
\[
e_r=(r,r),\qquad r\in S.
\]
Moreover,
\[
F_r=\{x\in M:x^2=e_r\}.
\]
Indeed, \((r,r)^2=(r,r)\), while \((r,c)^2=(r,r)\ne(r,c)\) for \(c\ne r\).

\subsection{Automorphism group}

\begin{theoremstar}[2.1]
Every automorphism of the PAB magma is induced by a unique
permutation of the three types:
\[
\operatorname{Aut}(M,\cdot)=\{\widehat\rho:\rho\in S_3\},
\qquad
\widehat\rho(r,c)=(\rho(r),\rho(c)).
\]
In particular,
\[
\operatorname{Aut}(M,\cdot)\cong S_3.
\]

\end{theoremstar}
\begin{proof}
The three idempotents are exactly \(e_r=(r,r)\), so any automorphism
\(\phi\) permutes them. Write \(\phi(e_r)=e_{\rho(r)}\) for a unique
\(\rho\in S_3\). Since \(F_r=\{x\in M:x^2=e_r\}\), automorphisms preserve square
fibers, hence \(\phi(F_r)=F_{\rho(r)}\).

For distinct \(r,s\in S\),
\[
(r,s)=e_r\cdot e_s.
\]
Therefore
\[
\phi(r,s)=\phi(e_r\cdot e_s)=\phi(e_r)\cdot\phi(e_s)
=e_{\rho(r)}\cdot e_{\rho(s)}=(\rho(r),\rho(s)).
\]
The same formula holds on diagonal elements because
\(\phi(r,r)=e_{\rho(r)}=(\rho(r),\rho(r))\). Thus
\(\phi(r,c)=(\rho(r),\rho(c))\) for all \((r,c)\in M\). Conversely, every such
diagonal permutation preserves equality/inequality of rows and columns and the
complement relation on \(S\), so it is an automorphism.

\end{proof}
\subsection{Endomorphism monoid}

\begin{theoremstar}[2.2]
The endomorphism monoid of PAB is
\[
\operatorname{End}(M,\cdot)=\{\widehat\rho:\rho\in S_3\}\cup\{\varepsilon_0,\varepsilon_1,\varepsilon_2\},
\]
where \(\widehat\rho(r,c)=(\rho(r),\rho(c))\) and \(\varepsilon_t(x)=e_t=(t,t)\)
for all \(x\in M\). Thus \(|\operatorname{End}(M,\cdot)|=9\).

\end{theoremstar}
\begin{proof}
Let \(\phi:M\to M\) be an endomorphism. Since endomorphisms preserve
idempotents, each \(e_r=(r,r)\) is sent to a diagonal idempotent. Define
\(\rho:S\to S\) by \(\phi(e_r)=e_{\rho(r)}\). Because
\(F_r=\{x:x^2=e_r\}\), we have \(\phi(F_r)\subseteq F_{\rho(r)}\). There are three
cases.

If \(|\operatorname{Im}\rho|=3\), then \(\rho\) is a permutation. For \(r\ne s\),
\((r,s)=e_r\cdot e_s\), so
\[
\phi(r,s)=\phi(e_r\cdot e_s)=e_{\rho(r)}\cdot e_{\rho(s)}=(\rho(r),\rho(s)).
\]
The same formula holds on the diagonal, hence \(\phi=\widehat\rho\).

If \(|\operatorname{Im}\rho|=1\), then \(\rho(r)=t\) for all \(r\). For any
off-diagonal element \((r,s)=e_r\cdot e_s\),
\[
\phi(r,s)=\phi(e_r\cdot e_s)=e_t\cdot e_t=e_t.
\]
Also \(\phi(e_r)=e_t\) for all \(r\). Hence \(\phi=\varepsilon_t\), the constant
endomorphism with value \(e_t=(t,t)\).

If \(|\operatorname{Im}\rho|=2\), choose distinct \(a,b,c\in S\) with
\(\rho(a)=\rho(b)=p\), \(\rho(c)=q\), and \(p\ne q\). In the fiber \(F_a\), the
Steiner rule gives
\[
e_a\cdot(a,b)=(a,c).
\]
Applying \(\phi\),
\[
\phi(e_a)\cdot\phi(a,b)=\phi(a,c).
\]
Now \(\phi(e_a)=e_p\), and since \((a,b)=e_a\cdot e_b\), we have
\(\phi(a,b)=e_p\cdot e_p=e_p\). Thus the left-hand side is \(e_p\). On the other
hand,
\[
\phi(a,c)=\phi(e_a\cdot e_c)=e_p\cdot e_q=(p,q)\ne e_p.
\]
This contradiction rules out the two-valued case. The three cases exhaust all
maps \(S\to S\).

\end{proof}
\subsection{Row barrier}

\begin{propositionstar}[2.3]
For any \(x\in F_r\), \(y\in F_s\), with \(r\ne s\),
\[
\langle x,y\rangle=F_r\cup F_s.
\]
The third row is unreachable.

\end{propositionstar}
\begin{proof}
By AX0, every product of elements already in \(F_r\cup F_s\) remains in
one of those two rows. Hence \(\langle x,y\rangle\subseteq F_r\cup F_s\).

For the reverse inclusion, first generate \(F_r\). The generator \(x\) gives
\(x^2=e_r\). If \(x\ne e_r\), then \(x\) and \(e_r\) are two distinct elements of
\(F_r\), hence AX2 generates the third element of \(F_r\). If \(x=e_r\), then
\(x\cdot y=(r,s)\ne e_r\), so \(e_r\) and \((r,s)\) generate \(F_r\). Thus
\(F_r\subseteq\langle x,y\rangle\).

The same argument with \(y\) gives \(F_s\subseteq\langle x,y\rangle\): either
\(y\ne e_s\) and \(y,y^2=e_s\) generate \(F_s\), or \(y=e_s\) and
\(y\cdot x=(s,r)\ne e_s\) supplies a second element. Therefore
\(\langle x,y\rangle=F_r\cup F_s\).

\end{proof}
\subsection{Universal identities and row-confinement}

\begin{propositionstar}[2.4]
PAB satisfies the following identities for all \(x,y\in M\):
\[
(x\cdot y)(x\cdot y)=x\cdot x,
\]
\[
(x\cdot x)(y\cdot y)=(x\cdot y)(y\cdot x),
\]
\[
(x\cdot x)(y\cdot x)=(x\cdot y)(y\cdot y).
\]
It also satisfies, for all \(z,w\in M\),
\[
z\cdot(z\cdot w)=(z\cdot w)\cdot z.
\]

\end{propositionstar}
\begin{proof}
Write \(e_r=(r,r)\). Since \(x^2=e_{\operatorname{row}(x)}\), the first
identity follows from AX0:
\[
(x\cdot y)^2=e_{\operatorname{row}(x\cdot y)}=e_{\operatorname{row}(x)}=x^2.
\]
For the second identity, let \(r=\operatorname{row}(x)\) and
\(s=\operatorname{row}(y)\). If \(r=s\), then \(x\cdot y=y\cdot x\) inside the
Steiner fiber, and the second identity reduces to the first. If \(r\ne s\), then
\[
(x^2)(y^2)=e_r e_s=(r,s),
\]
while
\[
(xy)(yx)=(r,s)(s,r)=(r,s)
\]
by AX1. The third identity is similar. If \(r=s\), all factors lie in the same
commutative fiber. If \(r\ne s\), then
\[
(x^2)(yx)=e_r(s,r)=(r,s),
\qquad
(xy)(y^2)=(r,s)e_s=(r,s).
\]
Finally, put \(u=z\cdot w\). By AX0, \(u\) lies in the same row as \(z\). If
\(u=z\) the last identity is trivial; if \(u\ne z\), it follows from commutativity
of the Steiner operation inside that row.

For every triple \((x,y,z)\), both bracketings of the associator stay in the row
of \(x\):
\[
((x\cdot y)\cdot z)_{\mathrm{row}}=(x\cdot(y\cdot z))_{\mathrm{row}}=
\operatorname{row}(x).
\]
Thus non-associativity is an intra-row column phenomenon.

\end{proof}


\section{Elementary classification}

\subsection{Full cross-rule universe}

A cross-row input is a quadruple \((r_1,c_1,r_2,c_2)\in S^4\) with \(r_1\ne r_2\).
There are \(54\) such inputs. The diagonal \(\sigma\)-action is free on these
inputs, so they form \(18\) orbits of size \(3\). A \(\sigma\)-equivariant
cross-rule is determined by one output in \(S\) on each orbit. Hence
\[
|\operatorname{CrossRules}_\sigma|=3^{18}=387,420,489.
\]

\subsection{Six distinguished elementary cross-rules}

This paper restricts to the following six elementary rules:

\begin{center}
\small
\begin{longtable}{@{}>{\raggedright\arraybackslash}p{0.313\textwidth}>{\raggedright\arraybackslash}p{0.313\textwidth}>{\raggedright\arraybackslash}p{0.313\textwidth}@{}}
Rule & Output column & Name \\
\midrule
\(g_1\) & \(r_2\) & column-blind / PAB \\
\(g_2\) & \(c_2\) & transparent \\
\(g_3\) & \(c_1c_2\) if \(c_1\ne c_2\), else \(c_1\) & complement-transparent \\
\(g_4\) & \(c_1\) & echo \\
\(g_5\) & \(r_1\) & self-referential \\
\(g_6\) & \(r_2c_2\) if \(r_2\ne c_2\), else \(r_2\) & anti-complement \\
\end{longtable}
\end{center}

They are selected elementary formulas, not an exhaustive list of equivariant
cross-row rules.

\subsection[Diagonal maps and Omega-elem]{Diagonal maps and \(\Omega_{\mathrm{elem}}\)}

A \(\sigma\)-equivariant diagonal map is determined by three base values
\((d_0,d_1,d_2)\in S^3\), where \(d_k=\delta(0,k)\). Hence \(|\Delta_\sigma|=27\).
The standard diagonal is \(\delta_{\mathrm{std}}=(0,0,0)\), i.e. \(\delta(r,c)=r\).

Define
\[
\Omega_{\mathrm{elem}}=\{g_1,\ldots,g_6\}\times\Delta_\sigma.
\]
Then
\[
|\Omega_{\mathrm{elem}}|=6\cdot27=162.
\]

\subsection[Diagonal S3-orbit count]{Diagonal \(S_3\)-orbit count}

\begin{theoremstar}[3.1]
Under diagonal relabelling by \(S_3\), \(\Omega_{\mathrm{elem}}\)
decomposes into \(90\) orbits.

\end{theoremstar}
\begin{proof}
Burnside's lemma gives
\[
|\Omega_{\mathrm{elem}}/S_3|=\frac{1}{6}\sum_{\pi\in S_3}|\operatorname{Fix}(\pi)|.
\]
The fixed counts are
\[
|\operatorname{Fix}(\mathrm{id})|=162,
\quad
|\operatorname{Fix}(\sigma)|=162,
\quad
|\operatorname{Fix}(\sigma^2)|=162,
\]
and for each transposition \(\tau\),
\[
|\operatorname{Fix}(\tau)|=18.
\]
Therefore
\[
|\Omega_{\mathrm{elem}}/S_3|=\frac{162+162+162+18+18+18}{6}=90.
\]
The fixed counts are enumerated in Appendix A.

This theorem concerns the diagonal \(S_3\)-action only. It is not asserted to be a
classification by arbitrary magma isomorphism.

\end{proof}


\section[Associativity formula in Omega-elem]{Associativity formula in \(\Omega_{\mathrm{elem}}\)}

For any \((g,\delta)\in\Omega_{\mathrm{elem}}\), associativity decomposes by
row-pattern:
\[
\operatorname{Assoc}(g,\delta)=RRR+SRR+RRS+RSR+DIST.
\]

The row-pattern counts are:

\begin{center}
\small
\begin{longtable}{@{}>{\raggedright\arraybackslash}p{0.313\textwidth}>{\raggedright\arraybackslash}p{0.313\textwidth}>{\raggedright\arraybackslash}p{0.313\textwidth}@{}}
Pattern & Row condition & Count \\
\midrule
\(RRR\) & all three rows equal & 81 \\
\(SRR\) & first row distinct, second and third equal & 162 \\
\(RRS\) & first and second equal, third distinct & 162 \\
\(RSR\) & first and third equal, second distinct & 162 \\
\(DIST\) & all rows distinct & 162 \\
\end{longtable}
\end{center}

For \(\delta\) with base tuple \((d_0,d_1,d_2)\), define
\[
p_{\mathrm{dist}}=|\{d_0,d_1,d_2\}|,
\]
\[
d_s=|\{k\in S:d_k=k\}|,
\]
\[
\mathrm{fix}=|\{(r,c):\delta(r,c)=c\}|,
\]
and decompose
\[
\mathrm{fix}=\mathrm{fix}_{\mathrm{diag}}+\mathrm{fix}_{\mathrm{off}},
\]
where \(\mathrm{fix}_{\mathrm{diag}}\) counts fixed pairs \((r,c)\) with \(r=c\),
and \(\mathrm{fix}_{\mathrm{off}}\) counts fixed pairs with \(r\ne c\).

\subsection[RRR-formula]{\(RRR\)-formula}

The \(RRR\) contribution depends only on \(\delta\):

\begin{center}
\small
\begin{tabular}{@{}cccc@{}}
\(p_{\mathrm{dist}}\) & condition on \(d_s\) & \(RRR\) & fiber associativity \\
\midrule
1 & any & 57 & \(19/27\) \\
2 & 0 & 33 & \(11/27\) \\
2 & \(\ge1\) & 39 & \(13/27\) \\
3 & any possible & 27 & \(9/27\) \\
\end{tabular}
\end{center}

\subsection{Mixed-pattern formulas}

\begin{center}
\small
\begin{tabular}{@{}lcccc@{}}
Rule & \(RRS\) & \(RSR\) & \(SRR\) & \(DIST\) \\
\midrule
\(g_1\) & \(9\mathrm{fix}_{\mathrm{off}}\) & \(9\mathrm{fix}_{\mathrm{off}}\) & 162 & 0 \\
\(g_2\) & \(6\mathrm{fix}\) & \(6\mathrm{fix}\) & \(6\mathrm{fix}\) & 162 \\
\(g_3\) & \(72-2\mathrm{fix}\) & 54 & 54 & 54 \\
\(g_4\) & 162 & \(6\mathrm{fix}\) & 162 & 162 \\
\(g_5\) & \(18\mathrm{fix}_{\mathrm{diag}}\) & \(18\mathrm{fix}_{\mathrm{diag}}\) & 162 & 162 \\
\(g_6\) & \(6\mathrm{fix}\) & \(81-3\mathrm{fix}\) & \(6\mathrm{fix}\) & 54 \\
\end{tabular}
\end{center}

\subsection{Master formula theorem}

\begin{theoremstar}[4.1]
For every \((g,\delta)\in\Omega_{\mathrm{elem}}\), the number of
associative triples is determined by \(g\) and the diagonal parameters
\(p_{\mathrm{dist}},d_s,\mathrm{fix},\mathrm{fix}_{\mathrm{diag}},\mathrm{fix}_{\mathrm{off}}\), via the tables above.

For example,
\[
\operatorname{Assoc}(g_1,\delta)
=RRR(p_{\mathrm{dist}},d_s)+18\,\mathrm{fix}_{\mathrm{off}}+162.
\]

\end{theoremstar}
\begin{proof}
The five row-patterns partition \(M^3\), so associativity is the sum
of the corresponding pattern counts. If a triple has type \(RRR\), both
bracketings remain inside one row-fiber; the cross-rule is never invoked. Thus
\(RRR\) depends only on the one-row operation determined by \(\delta\). Since the
three row-fibers are \(\sigma\)-translates, it is enough to enumerate the 27
column triples in one fiber. Grouping these triples by equality pattern gives
exactly the \(RRR\)-table above.

For the mixed patterns \(SRR,RRS,RSR,DIST\), AX0 implies that both bracketings
always land in the row of the first argument; associativity is therefore a
column equality. Substituting each of the six elementary cross-rules reduces the
condition to either a constant truth value or to one of the fixed-point counts
\(\mathrm{fix},\mathrm{fix}_{\mathrm{diag}},\mathrm{fix}_{\mathrm{off}}\). This case analysis gives the mixed-pattern
table. The computational verification in Appendix A independently enumerates all
\(162\) elementary magmas and all \(729\) triples for each magma; it reports zero
mismatches against these formulas.

\end{proof}


\section{Elementary associativity spectrum}

\begin{theoremstar}[5.1 --- Elementary associativity spectrum]
Inside \(\Omega_{\mathrm{elem}}\), the associativity function takes exactly the following values:
\[
\begin{aligned}
\{&162,168,189,195,201,219,228,243,255,261,267,273,285,297,\\
&309,351,357,363,381,459,471,489,513,519,531,543,561,567\}.
\end{aligned}
\]
Thus the spectrum has \(28\) values. Every value is divisible by \(3\). The extremal values are
\[
\min_{\Omega_{\mathrm{elem}}}\operatorname{Assoc}=162,
\qquad
\max_{\Omega_{\mathrm{elem}}}\operatorname{Assoc}=567.
\]
Equivalently, after normalisation by \(|M|^3=729\),
\[
\frac{2}{9}\le
\frac{\operatorname{Assoc}(g,\delta)}{729}
\le
\frac{7}{9}
\qquad ((g,\delta)\in\Omega_{\mathrm{elem}}).
\]
The values divisible by \(27\) are exactly
\[
162,189,243,297,351,459,513,567.
\]

\end{theoremstar}
\begin{proof}
The master formula of Theorem 4.1 expresses associativity as a finite
function of one of the six elementary rules and the \(27\) diagonal maps. Applying
that formula to the \(162\) pairs \((g,\delta)\in\Omega_{\mathrm{elem}}\) gives the
listed set of \(28\) values. The exhaustive validation in Appendix A independently
enumerates all \(729\) triples for every object of \(\Omega_{\mathrm{elem}}\), agrees
with the formula in every case, and records witnesses for the extremal values.
One minimum witness is \((g_6,\delta=(1,2,0))\), with associativity \(162\); one
maximum witness is \((g_4,\delta=(0,1,2))\), with associativity \(567\). The
divisibility and anchor-value claims are read directly from the verified spectrum
table.

These spectrum statements are made only for \(\Omega_{\mathrm{elem}}\), not for the
full \(3^{18}\)-cross-rule universe.

\end{proof}


\section{Information-theoretic selection}

\subsection{Standard elementary selection}

For the six standard elementary magmas \((g_i,\delta_{\mathrm{std}})\), define
\[
\lambda(g_i)=\frac{\operatorname{Assoc}(g_i,\delta_{\mathrm{std}})}{729}.
\]
The conditional access information is
\[
H_{\mathrm{access}}(g)=I(C_{\mathrm{out}};C_1,C_2\mid R_1,R_2).
\]

\begin{center}
\small
\begin{tabular}{@{}lcccc@{}}
Rule & \(H_{\mathrm{storage}}\) & \(H_{\mathrm{access}}\) & Assoc & \(\lambda\) \\
\midrule
\(g_1\) & 0 & 0 & 219 & \(219/729\) \\
\(g_2\) & \(\log_2 3\) & \(\log_2 3\) & 273 & \(273/729\) \\
\(g_3\) & \(\log_2 3\) & \(\log_2 3\) & 285 & \(285/729\) \\
\(g_4\) & \(\log_2 3\) & \(\log_2 3\) & 561 & \(561/729\) \\
\(g_5\) & 0 & 0 & 489 & \(489/729\) \\
\(g_6\) & 0 & \(\log_2 3\) & 219 & \(219/729\) \\
\end{tabular}
\end{center}

Thus
\[
\min H_{\mathrm{access}}=\{g_1,g_5\},
\qquad
\min\lambda=\{g_1,g_6\},
\]
and
\[
\{g_1,g_5\}\cap\{g_1,g_6\}=\{g_1\}.
\]
This selects PAB.

\subsection{Parameter-free functional uniqueness}

For \(\beta>0\), define
\[
F_\beta(g)=H_{\mathrm{access}}(g)-\beta(1-\lambda(g))
=H_{\mathrm{access}}(g)+\beta(\lambda(g)-1).
\]

\begin{corollarystar}[6.1]
Among the six standard elementary magmas,
\(F_\beta\) is uniquely minimised by \(g_1\) for every \(\beta>0\).

\end{corollarystar}
\begin{proof}
Against \(g_5\), the access costs are equal, while
\[
\lambda(g_1)-\lambda(g_5)=\frac{219-489}{729}<0.
\]
Therefore \(F_\beta(g_1)-F_\beta(g_5)<0\). For
\(g\in\{g_2,g_3,g_4,g_6\}\),
\[
H_{\mathrm{access}}(g_1)-H_{\mathrm{access}}(g)=-\log_2 3<0,
\]
and
\[
\lambda(g_1)-\lambda(g)\le0.
\]
Thus
\[
F_\beta(g_1)-F_\beta(g)
=-\log_2 3+
\beta(\lambda(g_1)-\lambda(g))
\le -\log_2 3<0.
\]
So \(F_\beta(g_1)<F_\beta(g)\) for every \(g\ne g_1\).

\end{proof}
\subsection[Restricted sequential selection in Omega-prime-elem]{Restricted sequential selection in \(\Omega'_{\mathrm{elem}}\)}

Let \(\mathcal F_\sigma\) be the nine \(\sigma\)-equivariant distinct-column within-fiber rules and
\[
\Omega'_{\mathrm{elem}}
=\mathcal F_\sigma\times\{g_1,\ldots,g_6\}\times\Delta_\sigma.
\]
Then \(|\Omega'_{\mathrm{elem}}|=1458\).

The verified sequential chain is:

\begin{center}
\small
\begin{longtable}{@{}>{\raggedright\arraybackslash}p{0.313\textwidth}>{\raggedright\arraybackslash}p{0.313\textwidth}>{\raggedright\arraybackslash}p{0.313\textwidth}@{}}
Stage & Criterion & Remaining candidates \\
\midrule
0 & full \(\Omega'_{\mathrm{elem}}\) & 1458 \\
1 & minimise \(H_{\mathrm{access}}\) & 486 \\
2 & minimise \(H_{\mathrm{diag}}\) among Stage 1 & 54 \\
3 & minimise associativity among Stage 2 & 1 \\
\end{longtable}
\end{center}

The unique final remaining candidate is
\[
(g_1,F_{21},\delta_{\mathrm{std}})=PAB.
\]
Here \(F_{21}\) is the complement / Steiner fiber rule.

The statement is sequential: the associativity criterion is applied after the two
entropy-based filters. Over all of \(\Omega'_{\mathrm{elem}}\), the global
associativity minimum is \(135\), not \(219\), and is not attained by PAB.



\section{Absorption graph}

Let
\[
M^\times=\{(r,c)\in M:r\ne c\}.
\]
For \(x=(r,c)\in M^\times\), define
\[
T_\alpha(r,c)=(c,rc),
\qquad
T_\beta(r,c)=(c,r).
\]
For \(x\in M^\times\), define
\[
\operatorname{Abs}(x)=\{y\in M^\times\setminus\{x\}:x\cdot y=x\}.
\]

\begin{theoremstar}[7.1]
For every \(x\in M^\times\),
\[
\operatorname{Abs}(x)=\{T_\alpha(x),T_\beta(x)\}.
\]

\end{theoremstar}
\begin{proof}
Let \(x=(r,c)\in M^\times\), so \(r\ne c\). If \(y=(s,d)\in M^\times\)
has \(s=r\), then \(x\cdot y\) is computed inside \(F_r\). If \(y\ne x\), the
Steiner rule gives the third element of the fiber, not \(x\). Hence no same-row
off-diagonal element is absorbed by \(x\).

If \(s\ne r\), then AX1 gives \(x\cdot y=(r,s)\). This equals \(x=(r,c)\) if and
only if \(s=c\). The off-diagonal elements with row \(c\) are exactly
\[
(c,r)=T_\beta(x)
\]
and
\[
(c,rc)=T_\alpha(x).
\]
Therefore \(\operatorname{Abs}(x)=\{T_\alpha(x),T_\beta(x)\}\).

\end{proof}


\section{Minimal finite-geometric interpretation}

The set \(M=S\times S\) is the affine plane \(AG(2,3)\). Each row fiber is a
three-point line, and AX2 is the corresponding Steiner completion rule on that
line. The six off-diagonal states
\[
M^\times=\{(r,c):r\ne c\}
\]
form the usual hexagonal \(A_2\)-root pattern. The maps \(T_\alpha\) and
\(T_\beta\) generate the derived absorption structure on this hexagon.

\appendix
\section{Computational verification}

\subsection{Verification policy}

All finite claims in the paper are verified by exact exhaustive enumeration over
finite sets. The verification scripts use integer arithmetic and symbolic labels
in \(S=\{0,1,2\}\); no floating-point comparison is used for algebraic claims.

All scripts and generated artifacts are available in the public repository:
\begin{center}
\texttt{https://github.com/burundai-t/sigma-equivariant-magmas-ag23}
\end{center}
The repository mirrors the directory structure used in this appendix and allows full reproduction of all finite verification results.

Each verification script:
\begin{enumerate}[leftmargin=2.2em,itemsep=0.25em]
\item defines the finite state space explicitly;
\item enumerates the relevant objects exhaustively;
\item writes stable CSV/log outputs;
\item exits with failure if any required assertion is violated.
\end{enumerate}

\subsection{Finite state spaces}

\begin{center}
\small
\begin{tabular}{@{}ll@{}}
\toprule
Object & Size \\
\midrule
\(M=S\times S\) & 9 \\
\(M^2\) & 81 \\
\(M^3\) & 729 \\
Cross-row inputs \((r_1,c_1,r_2,c_2)\), \(r_1\ne r_2\) & 54 \\
\(\sigma\)-orbits of cross-row inputs & 18 \\
Full \(\sigma\)-equivariant cross-rules & \(3^{18}=387{,}420{,}489\) \\
Elementary cross-rules & 6 \\
\(\sigma\)-equivariant diagonal maps & 27 \\
Elementary classification space \(\Omega_{\mathrm{elem}}\) & 162 \\
Extended elementary space \(\Omega'_{\mathrm{elem}}\) & 1458 \\
\bottomrule
\end{tabular}
\end{center}

\subsection{Reproducibility inventory}

The artifact root is \texttt{proof\_artifacts/}. Each block below corresponds to a repository subdirectory; the table lists the script, verified claim, and primary outputs without repeating the full path.

\begin{center}
\scriptsize
\setlength{\tabcolsep}{4pt}
\begin{longtable}{@{}p{0.10\textwidth}p{0.22\textwidth}p{0.27\textwidth}p{0.31\textwidth}@{}}
\toprule
Block & Script & Verified claim & Primary outputs \\
\midrule
A1 &
{\ttfamily cross\_rule\_count.py} &
Full \(\sigma\)-equivariant cross-rule universe has cardinality \(3^{18}\); the six \(g_i\) are distinguished elementary rules. &
{\ttfamily cross\_rule\_orbits.csv}\newline
{\ttfamily elementary\_rule\_}\newline
{\ttfamily vectors.csv} \\
\addlinespace
A2 &
{\ttfamily a2\_burnside\_}\newline
{\ttfamily fixed\_table.py} &
\(\Omega_{\mathrm{elem}}\) has 90 diagonal \(S_3\)-orbits under relabelling. &
{\ttfamily A2\_burnside\_}\newline
{\ttfamily fixed\_table.csv}\newline
{\ttfamily A2\_burnside\_orbits.csv} \\
\addlinespace
A3 &
{\ttfamily a3\_mixed\_pattern\_}\newline
{\ttfamily formula\_table.py} &
Master formula for associativity over \(\Omega_{\mathrm{elem}}\), decomposed by row-pattern. &
{\ttfamily mixed\_pattern\_}\newline
{\ttfamily formula\_table.csv}\newline
{\ttfamily rrr\_formula\_table.csv}\newline
{\ttfamily mixed\_pattern\_}\newline
{\ttfamily formula\_validation.csv} \\
\addlinespace
A4 &
{\ttfamily a4\_spectrum\_}\newline
{\ttfamily values.py} &
Elementary associativity spectrum has exactly 28 values in the range 162--567. &
{\ttfamily spectrum\_values.csv}\newline
{\ttfamily spectrum\_}\newline
{\ttfamily witnesses.csv}\newline
{\ttfamily spectrum\_anchor\_}\newline
{\ttfamily values.csv} \\
\addlinespace
A5 &
{\ttfamily omega\_prime\_}\newline
{\ttfamily selection.py} &
Restricted sequential selection in \(\Omega'_{\mathrm{elem}}\): \(1458\to486\to54\to1\). &
{\ttfamily omega\_prime\_}\newline
{\ttfamily fiber\_rules.csv}\newline
{\ttfamily omega\_prime\_survivors.csv}\newline
{\ttfamily omega\_prime\_}\newline
{\ttfamily selection\_steps.csv}\newline
{\ttfamily omega\_prime\_}\newline
{\ttfamily step3\_winner.csv}\newline
{\ttfamily omega\_prime\_}\newline
{\ttfamily global\_min\_assoc.csv} \\
\addlinespace
S1 &
{\ttfamily p1\_polishing\_}\newline
{\ttfamily verification.py} &
Automorphism rigidity, endomorphism monoid, standard elementary selection, symbolic \(F_\beta\), and absorption graph checks. &
{\ttfamily automorphism\_table.csv}\newline
{\ttfamily endomorphism\_table.csv}\newline
{\ttfamily standard\_canonical\_}\newline
{\ttfamily selection.csv}\newline
{\ttfamily fbeta\_symbolic\_}\newline
{\ttfamily comparisons.csv}\newline
{\ttfamily absorption\_graph\_}\newline
{\ttfamily table.csv} \\
\bottomrule
\end{longtable}
\end{center}

\subsection{Corrected cross-rule count}

The cross-row input set is
\[
\Sigma=\{(r_1,c_1,r_2,c_2)\in S^4:r_1\ne r_2\},
\qquad |\Sigma|=54.
\]
The diagonal cyclic action of \(\sigma\) on \(\Sigma\) is free, so it has 18
orbits. A \(\sigma\)-equivariant cross-rule is determined by one value in \(S\)
on each orbit, hence
\[
|\mathrm{CrossRules}_\sigma|=3^{18}=387{,}420{,}489.
\]
The corresponding verification block is \texttt{A1\_cross\_rule\_count}. Thus
the six rules \(g_1,\ldots,g_6\) are distinguished elementary cross-rules, not
the full set of \(\sigma\)-equivariant cross-rules.

\subsection{Burnside count for \texorpdfstring{\(\Omega_{\mathrm{elem}}\)}{Omega-elem}}

For
\[
\Omega_{\mathrm{elem}}=\{g_1,\ldots,g_6\}\times\Delta_\sigma,
\qquad |\Omega_{\mathrm{elem}}|=162,
\]
the diagonal \(S_3\)-relabelling fixed counts are:

\begin{center}
\small
\begin{tabular}{@{}lr@{}}
\toprule
Group element & Fixed objects \\
\midrule
identity & 162 \\
3-cycle \(\sigma\) & 162 \\
3-cycle \(\sigma^2\) & 162 \\
transposition \(\tau_{01}\) & 18 \\
transposition \(\tau_{02}\) & 18 \\
transposition \(\tau_{12}\) & 18 \\
\bottomrule
\end{tabular}
\end{center}

Thus Burnside gives
\[
\frac{162+162+162+18+18+18}{6}=90.
\]
The verified statement is that \(\Omega_{\mathrm{elem}}\) decomposes into 90
diagonal \(S_3\)-orbits. These are not asserted to be full magma isomorphism
classes.

\subsection{Associativity master formula}

For each \((g,\delta)\in\Omega_{\mathrm{elem}}\), the computation enumerates all
729 triples and decomposes associativity into row-patterns
\[
RRR,\quad SRR,\quad RRS,\quad RSR,\quad DIST.
\]
The formula table is validated for all 162 objects. The validation output includes
\begin{verbatim}
Formula mismatches: 0
\end{verbatim}
The corresponding verification block is \texttt{A3\_mixed\_pattern\_formula}.

\subsection{Associativity spectrum in \texorpdfstring{\(\Omega_{\mathrm{elem}}\)}{Omega-elem}}

The verified spectrum is
\[
\begin{gathered}
162,168,189,195,201,219,228,243,255,261,267,273,285,297,\\
309,351,357,363,381,459,471,489,513,519,531,543,561,567.
\end{gathered}
\]
It contains exactly 28 values. The minimum is 162, the maximum is 567, and all
values are divisible by 3. The values divisible by 27 are
\[
162,189,243,297,351,459,513,567.
\]
The corresponding verification block is \texttt{A4\_spectrum}.

\subsection{Standard elementary selection}

The standard elementary selection table is:

\begin{center}
\small
\begin{tabular}{@{}lllll@{}}
\toprule
Rule & \(H_{\mathrm{storage}}\) & \(H_{\mathrm{access}}\) & Assoc at \(\delta_{\mathrm{std}}\) & Selection role \\
\midrule
\(g_1\) & 0 & 0 & 219 & selected \\
\(g_2\) & \(\log_2 3\) & \(\log_2 3\) & 273 & rejected \\
\(g_3\) & \(\log_2 3\) & \(\log_2 3\) & 285 & rejected \\
\(g_4\) & \(\log_2 3\) & \(\log_2 3\) & 561 & rejected \\
\(g_5\) & 0 & 0 & 489 & rejected by associativity \\
\(g_6\) & 0 & \(\log_2 3\) & 219 & rejected by access cost \\
\bottomrule
\end{tabular}
\end{center}

Thus
\[
\min H_{\mathrm{access}}=\{g_1,g_5\},
\qquad
\min\lambda=\{g_1,g_6\},
\]
and the intersection is \(\{g_1\}\). The normalized table and the symbolic \(F_\beta\)-comparisons are recorded in the structural verification outputs listed above.

\subsection{Restricted sequential selection in \texorpdfstring{\(\Omega'_{\mathrm{elem}}\)}{Omega-prime-elem}}

The extended elementary space is
\[
\Omega'_{\mathrm{elem}}
=\mathcal F_\sigma\times\{g_1,\ldots,g_6\}\times\Delta_\sigma,
\qquad |\Omega'_{\mathrm{elem}}|=1458.
\]
The nine fiber rules are parametrized by seeds
\[
u=f(0,1),\qquad v=f(0,2),
\]
with equivariance determining all other distinct-column products. The
PAB/Steiner rule is \(F_{21}\).

The verified sequential chain is:
\begin{center}
\small
\begin{tabular}{@{}lll@{}}
\toprule
Stage & Criterion & Remaining candidates \\
\midrule
0 & full \(\Omega'_{\mathrm{elem}}\) & 1458 \\
1 & minimise \(H_{\mathrm{access}}\) & 486 \\
2 & minimise \(H_{\mathrm{diag}}\) among Stage 1 & 54 \\
3 & minimise associativity among Stage 2 & 1 \\
\bottomrule
\end{tabular}
\end{center}
The unique final remaining candidate is
\[
(g_1,F_{21},\delta_{\mathrm{std}})=\mathrm{PAB}.
\]
The verification also records that the global associativity minimum over all
\(\Omega'_{\mathrm{elem}}\) is 135 and is not attained by PAB. Therefore the
selection theorem is stated sequentially.

\subsection{Additional structural checks for PAB}

The structural verification block checks the finite parts of the PAB-specific
claims used in the main text. It enumerates the six diagonal automorphisms, the
six automorphisms plus three constant endomorphisms, the standard elementary
selection table, the symbolic \(F_\beta\)-comparisons, and the six-row
absorption graph table.

\section{Open problems and further directions}
\begin{enumerate}[leftmargin=2.2em,itemsep=0.35em]
\item Determine whether the 90 diagonal \(S_3\)-orbits in
\(\Omega_{\mathrm{elem}}\) coincide with full magma isomorphism classes.
Equivalently, prove or refute the required rigidity lemma.
\item Classify or stratify the full \(3^{18}\)-element \(\sigma\)-equivariant
cross-rule universe.
\item Determine whether the free PAB-magma on two generators is finite or
infinite.
\item Define and analyze operator-level path-memory invariants for PAB and for
restricted elementary families.
\item Develop the operator linearization theory of PAB, including left-multiplication
operators and intertwining spaces.
\end{enumerate}

\section*{Note on AI assistance}

This work was developed with the assistance of large language models.

The author specified the conceptual framework, constraints, and target properties 
of the objects under study. Within this framework, AI tools were used for 
computational exploration, code drafting, and structuring intermediate results.

All mathematical definitions, statements, and theorems presented in the paper 
were established and independently verified by the author, with exhaustive 
finite checks (Appendix~A).

\bibliographystyle{alpha}
\bibliography{references}
\end{document}
